\documentclass[12pt,letterpaper]{article}

% --- Paquetes ---
\usepackage[utf8]{inputenc}
\usepackage[T1]{fontenc}
\usepackage[spanish]{babel}
\usepackage[margin=2.5cm]{geometry}
\usepackage{float}
\usepackage{tabularx}
\usepackage{hyperref}
\usepackage{graphicx}
\usepackage{setspace}
\usepackage{parskip}
\usepackage{helvet}      % Fuente sans-serif para un look más limpio
\renewcommand{\familydefault}{\sfdefault}

% --- Datos de la portada ---
\newcommand{\institucion}{Universidad Nacional Autónoma de México}
\newcommand{\facultad}{Facultad de Ingeniería}
\newcommand{\numpractica}{Práctica 3}
\newcommand{\semestre}{Semestre 2027-1}
\newcommand{\entrega}{24 de septiembre de 2026}
\newcommand{\grupo}{Grupo: 16}
\newcommand{\alumno}{López Guevara Sebastián}
\newcommand{\profesora}{Sonia López Maldonado}
\newcommand{\practica}{Ciclo Rankine}
\newcommand{\materia}{SISTEMAS DE CONVERSIÓN DE ENERGÍA TÉRMICA}

\begin{document}
\thispagestyle{empty}

\begin{center}

    % --- INSTITUCIÓN ---
    {\Large\bfseries \institucion}\\[0.3cm]
    {\Large\bfseries \facultad}\\[0.3cm]
    {\large \semestre}\\[0.2cm]
    {\large \grupo}\\[1.5cm]
    
    % --- LOGO (opcional) ---
    % \includegraphics[width=0.25\textwidth]{logo.png}\\[1.5cm]
    
    % --- TÍTULO DE LA PRÁCTICA ---
    {\large \numpractica}\\[0.2cm]
    {\Huge\bfseries \practica}\\[0.5cm]
    \rule{0.8\textwidth}{0.4pt}\\[1.5cm]

    % --- DATOS DEL ALUMNO Y PROFESORA ---
    \begin{tabular}{rl}
        \textbf{Alumno:}     & \alumno     \\[0.3cm]
        \textbf{Profesora:}  & \profesora  \\[0.3cm]
        \textbf{Materia:}    & \materia    \\
    \end{tabular}

    \vfill
    
    % --- FECHA ---
    {\entrega}
    
\end{center}

\break

\section{Objetivos}

Obtener numericamente a partir de mediciones realizadas a un generador
de ciclo rankine en operacion las eficiencias interna, mecanica y total de la turbina, y las
eficiencias teorica y termica del ciclo completo. Esto utilizando los datos ya calculados 
para medir la calidad del vapor a la salida de la caldera.

\section{Materiales}

\begin{itemize}
    \item 1 Termometro 150 $^\circ C$
    \item 1 Termometro 100 $^\circ C$
    \item 1 Tacometro
    \item 1 Cronometro
    \item 3 Cubetas
\end{itemize}

\section{Desarrollo}

\begin{enumerate}
    % Etapa: 
    \item Identificar manometros y pozos para termometros
    % Etapa: 
    \item Identificar medidor de corriente y voltaje
    \item Primero se dejo pasar vapor hacia la turbina y se fue calibrando con el reostato
    \item Al llegar aproximadamente a los 120 volts, se tomo la medida de la corriente y del voltaje
    \item Al mismo tiempo, se recolecto el agua que salia del pozo para evitar que se tirara
    \item Se dejo llenar una cubeta durante 1 minuto. Se peso y luego se vacio para 
    obtener el peso del condensado. 

% Presiones

    % Etapa: 1 -> Entrada a la valvula | Salida del tanque de condensado (cabezal)
    \item Se tomo la lectura del manometro de la valvula 
    % Etapa: 2 -> Entrada a la turbina
    \item Se tomo la lectura del manometro a la entrada de la turbina
    % Etapa: 3 -> Salida de la turbina
    \item Se tomo la lectura del manometro a la salida de la turbina

% Temperaturas

    % Etapa: 3 -> Salida de la turbina
    \item Se tomo la lectura del termometro a la salida de la turbina
    % Etapa: 4 -> Salida del condensado
    \item Se tomo la lectura del termometro a la salida del condensador
    % Etapa: 5 -> Temperatura del tanque del condensado
    \item Se tomo la lectura del termometro dentro del tanque de condensado

\end{enumerate}

\section{Tabla de datos}

\begin{table}[H] \label{tab:data}
\centering
\renewcommand{\arraystretch}{1.4}
    \begin{tabular}{|l|l|}
    \hline
    \textbf{Medición} & \textbf{Variable} \\
    \hline
    $11.0\ [V]$  & $V$ \\
    \hline
    $61.3\ [A]$  & $A$ \\
    \hline
    $6.5\ \left[\frac{kg}{cm^2}\right]$  
                 & $P_1$ \\
    \hline
    $6\ \left[\frac{kg}{cm^2}\right]$  
                 & $P_2$ \\
    \hline
    $20\ [cm Hg]$  & $P_3$ \\
    \hline
    $73\ [^\circ C]$  & $T_3$ \\
    \hline
    $31\ [^\circ C]$  & $T_4$ \\
    \hline
    $36\ [^\circ C]$  & $T_5$ \\
    \hline
    $5.5\ [kg]$  & $M_{cond}$ \\
    \hline
    $300\ [g]$  & $M_{cub}$ \\
    \hline
    $1\ [min]$  & $t$ \\
    \hline
    \end{tabular}
\caption{Datos experimentales obtenidos durante la práctica.}
\end{table}

\section{Cálculos}

A continuacion se muestran los calculos realizados para obtener las metricas propuestas en
los objetivos.

Los valores encontrados en tablas se obtuvieron a partir de un script en python con
ayuda de la paqueteria CoolProp. Esto favorece la precision y rapidez para obtener dichos datos.

\subsection{Obtencion de las propiedades del estado I} \label{sec:estadoI}

Establecemos a la presion atmosferica con un valor de 
$78.51[kPa]$. Y utilizando los valos pertenecientes al estado 1 de la tabla de datos,
obtenemos la presion absoluta.

Realizando la conversion a $[Pa]$

$$P_1 = 6.5\left[\frac{kg}{cm^2}\right] 
\cdot 
\frac{98066.5[Pa]}{1\left[\frac{kg}{cm^2}\right]}=637432.25[Pa]=637.43225[kPa]
$$

$$P_{1_{abs}} = 637.43225[kPa] + 78.51[kPa] = 715.94225[kPa]$$

Y para entrar a tablas, tenemos el dato de la calidad del vapor a la salida de la caldera ($x=0.97$).

La entalpia del estado 1 es
$$
    h_1 =  2701.8194 \left[\frac{kJ}{kg}\right]
$$

\subsection{Obtencion de las propiedades del estado II} \label{sec:estadoII}

Despues de que el vapor sale de la caldera, y pasa por el cabezal, llega a la valvula "de paso?".

Al ser un proceso isentalpico, le asignamos la misma entalpia que a nuestro estado 1.

$$
    h_2 =  2701.8194 \left[\frac{kJ}{kg}\right]
$$

Para el siguiente estado necesitamos de la entropia. Para poder obtener este valor de tablas, 
convertimos la presion del estado 2 a $[kPa]$.

$$P_2 = 6\left[\frac{kg}{cm^2}\right] 
\cdot 
\frac{98066.5[Pa]}{1\left[\frac{kg}{cm^2}\right]}=588399.0[Pa]=588.399[kPa]
$$

Obtenemos la presion absoluta

$$P_{2_{abs}} = 588.399[kPa] + 78.51[kPa] = 666.909[kPa]$$

y entrando a tablas con los valores de entalpia y presion obtenemos

$$
    s_2 = 6.58849418722697  \left[\frac{kJ}{kg\ K}\right]
$$

\subsection{Obtencion de las propiedades del estado III} \label{sec:estadoIII}

Sabemos que el proceso termodinamico que ocurre en la turbina es isentropico idealmente.
Para este primer calculo de $h_{3_{isen}}$, 

\subsubsection{Calculo de la entalpia teorica}

Igualamos entropias

$$
    s_3 = s_2
$$

y convertimos el dato de presion vacuometrica del estado 3 a presion absoluta

$$
    P_3 = 20[cm\ Hg] \cdot \frac{1322.2[Pa]}{1[cm\ Hg]} = 26664.3913[Pa] = 26.6643913[kPa] 
$$

$$
    P_{3_{abs}} = 78.51[kPa] - 26.6643913[kPa] = 51.8456[kPa]
$$

Con estos dos datos ($s_3$,$P_{3_{abs}}$) entramos a tablas para obtener la entalpia

$$
    h_{3_{isen}} = 2294.1257 \left[\frac{kJ}{kg}\right]
$$

\subsubsection{Calculo de la entalpia real}

Para obtener la entalpia real utilizamos la presion absoluta del estado 3 ya obtenida y 
la temperatura medida a la salida de la turbina. Con

$$
    P_{3_{abs}} = 78.51[kPa] - 26.6643913[kPa] = 51.8456[kPa]
$$

$$
    T_3 = 73[^\circ C]
$$

obtenemos una entalpia real de

$$
    h_{3_{real}} = 305.6553 \left[\frac{kJ}{kg}\right]
$$

\subsection{Obtencion de las propiedades del estado IV} \label{sec:estadoIV}

A la entrada del condensador medimos la temperatura. Para calcular la entalpia utilizamos

\begin{equation}
    h = C_p\ T
\end{equation}

sustituyendo

$$
    h_4 = 4.186\left[\frac{kJ}{kg\ ^\circ C}\right] \cdot 31[^\circ C]
    = 0.129766 \left[\frac{kJ}{kg}\right]
$$

\subsection{Obtencion de las propiedades del estado V}

De la misma forma, conocemos la temperatura dentro del tanque de condensado, 
y sustituyendo en la ecuacion anterior

$$
    h_5 = 4.186\left[\frac{kJ}{kg\ ^\circ C}\right] \cdot 36[^\circ C]
    = 0.150696 \left[\frac{kJ}{kg}\right]
$$

\subsection{Obtencion de las eficiencias en la turbina}

\subsubsection{Eficiencia interna}

Para encontrar la eficiencia interna de la turbia utilizaremos las entalpias de los estados 3 real, 3 teorico y 2.

Dada la ecuacion 

\begin{equation}
    \eta_{int} = \frac{h_2 - h_{3_R}}{h_2 - h_{3_T}}
\end{equation}

Sustituimos los datos de los estados ya mencionados encontrados en 
las secciones \ref{sec:estadoII} y \ref{sec:estadoIII}.

Sustituimos

$$
    \eta_{int} = \frac{2701.8194 [kJ/kg] - 305.65531 [kJ/kg]}
    {2701.8194 [kJ/kg] - 2294.1257 [kJ/kg]} = 5.8773
$$



\subsubsection{Eficiencia mecanica}

Para encontrar la eficiencia mecanica de la turbina utilizamos

\begin{equation} \label{eq:efMec}
    \eta_{mec} = \frac{VI}{G_{v_{max}}(h_2 - h_{3_R})}
\end{equation}

Para poder evaluar necesitamos obtener GvMax. 

\begin{equation}
    G_{v_{max}} = \frac{m_{con} - m_{cu}}{t}    
\end{equation}

Utilizando los datos de tablas sustituimos y evaluamos 

$$
    G_{v_{max}} = \frac{5.5[kg] - 0.3[kg]}{1[min]} \cdot \frac{60[min]}{1[h]} = 312 \left[\frac{kg}{h}\right]
$$

Con este gasto, sustituimos en la ecuacion \ref{eq:efMec} con los datos obtenidos en 
las secciones \ref{sec:estadoII} y \ref{sec:estadoIII}.

$$
    \eta_{mec} = \frac{
        11[V]\cdot 61.3[A]
    }{
        312 \left[\frac{kg}{h}\right] 
        ( 2701.8194 \left[\frac{kJ}{kg}\right] - 305.65531 \left[\frac{kJ}{kg}\right] )
    }
    = 
    0.0009
$$

\subsection{Obtencion de las eficiencias del ciclo}

Para obtener la eficiencia del ciclo podemos sustituir directamente, ya que tenemos todos los datos.

Para la eficiencia teorica

\begin{equation}
    \eta_{teorica} = \frac{h_2 - h_{3_T}}{h_2 - h_5} 
\end{equation}

Sustituyendo 

$$
    \eta_{teorica} = \frac{
        2701.8194 \left[\frac{kJ}{kg}\right] - 2294.1257 \left[\frac{kJ}{kg}\right]
    }{
        2701.8194 \left[\frac{kJ}{kg}\right] - 0.1506 \left[\frac{kJ}{kg}\right]
    }
$$

\begin{equation}
    \eta_{termica} = \frac{h_2 - h_{3_R}}{h_2 - h_5} 
\end{equation}

Sustituyendo

$$
    \frac{
        2701.8194 \left[\frac{kJ}{kg}\right] - 305.6553 \left[\frac{kJ}{kg}\right]
    }{
        2701.8194 \left[\frac{kJ}{kg}\right] - 0.1506 \left[\frac{kJ}{kg}\right]
    }
$$

\section{Tabla de resultados}
\section{Analisis de resultados}
\section{Bibliografia}

% "https://www.google.com/url?sa=t&source=web&rct=j&opi=89978449&url=http://189.254.33.151/stn/coyoacan/Current_Monitor.htm&ved=2ahUKEwj9g9GOzfiWAxWHnCYFHfYLAX8QFnoECEAQAQ&usg=AOvVaw1rIQ2RuxiJD75yRA1eR6PM"

\end{document}